\documentclass[conference]{IEEEtran}
\IEEEoverridecommandlockouts

\usepackage[utf8]{inputenc}
\usepackage{graphicx}
\usepackage{cite}
\usepackage[hidelinks]{hyperref}

\begin{document}

\title{Comics and Scale in \textit{Watchmen} Chapter 1 and Social Media Comics:\\
How Form Reframes Power, Intimacy, and the Everyday}

\author{\IEEEauthorblockN{Ma Toan Bach}
\IEEEauthorblockA{\textit{Seneca Polytechnic}\\
Toronto, Canada \\
mbach@myseneca.ca}
}

\maketitle

\section{Introduction}
Comics are built to manipulate scale: not just the literal size of figures in a panel, but the perceived scale of events, emotions, and ideas. Scott McCloud argues that comics create meaning through the gutter and closure---the reader’s act of mentally completing what happens between panels \cite{McCloud1993}. He also shows that time and pacing are controlled by panel arrangement, and that word--image relationships can reinforce or complicate what we see \cite{McCloud1993}. Using these tools, \textit{Watchmen} Chapter 1 and short social media comics (CatanaComics and \textit{Strange Planet}) both ``scale'' experience up or down, but toward different ends: \textit{Watchmen} uses scale to dramatize vulnerability within political power, while social media comics use scale to inflate everyday intimacy and absurdity into something emotionally or conceptually ``big'' \cite{MooreGibbonsWatchmen,CourseWeek12,CatanaComics,Pyle2019}.

\section{Watchmen and the Politics of Scale}
\textit{Watchmen} opens with one of the clearest demonstrations of scale as a storytelling device: a tight close-up of the blood-smeared smiley-face badge that gradually expands into a broader view of the street and surrounding city \cite{MooreGibbonsWatchmen}. The shift from micro (a small, cheerful symbol) to macro (the bleak urban scene that contains it) makes the badge feel like more than evidence---it becomes a thematic anchor. The murder itself is not shown directly; the comic relies on closure as the reader infers what happened from aftermath details and sequential clues. That gap matters: by withholding the moment of impact, the comic lets the unseen act loom larger than any single image could. In other words, \textit{Watchmen} uses the gutter to enlarge the moral and political weight of violence. The small object in the foreground (the badge) becomes a portal into a much larger atmosphere of dread and corruption.

% ---------- Optional Figure Placeholder (insert your Watchmen badge/zoom-out image) ----------
% \begin{figure}[!t]
%     \centering
%     \includegraphics[width=\linewidth]{watchmen-zoomout-placeholder.png}
%     \caption{\textit{Watchmen} Ch.~1: close-up-to-city ``zoom-out'' sequence (insert your screenshot).}
%     \label{fig:watchmen-zoomout}
% \end{figure}
% -------------------------------------------------------------------------------

This sense of scale is reinforced by pacing. Many pages in Chapter 1 use a strict, repeating grid that produces a controlled rhythm, like a visual metronome \cite{MooreGibbonsWatchmen}. That regularity can ``box in'' characters and events, suggesting a world governed by systems---law, surveillance, ideology---rather than individual freedom. When \textit{Watchmen} does change scale, the change becomes more dramatic because it breaks the expected tempo. A striking example is the scene where Rorschach confronts Dr. Manhattan and Laurie: Manhattan’s presence is emphasized through panels that visually prioritize him, making him seem literally larger-than-life \cite{MooreGibbonsWatchmen}. The effect is not only physical; it’s temporal. Larger or more dominant panels slow the reader down, forcing attention, while smaller conversational panels speed up the beat of human dialogue. The sequence makes Manhattan feel as if he exists at a different ``size'' of reality---less bound to ordinary time, less emotionally available---while the humans remain confined to quick, limited exchanges. Here, panel scale becomes power scale: Manhattan’s bigness equals authority and detachment, while Rorschach and Laurie read as comparatively small and vulnerable within the same layout.

The word--image relationship in \textit{Watchmen} also manipulates scale. Rorschach’s narration and the visual field do not simply match; they often create tension between an expansive moral voice and a grim, concrete world \cite{MooreGibbonsWatchmen}. Even when the words sound absolute, the images undercut certainty with ambiguity and ugliness. This contrast magnifies the political atmosphere: the comic feels ``large'' not because it shows armies or global maps, but because it frames individual moments as symptoms of wider structures---Cold War paranoia, public fear, and institutional power. \textit{Watchmen} scales up the personal until it feels historical.

\section{Social Media Comics and Platform Scale}
Social media comics work differently, but they also depend on the same core mechanics. A CatanaComics post in the Week 12 set uses four panels to turn a tiny event---shopping for a used item online---into a joke about modern life and relationships \cite{CourseWeek12,CatanaComics}. The first panel sets up a small-scale question (``you got the goods?'' / ``you got the \$7?''), then the strip jumps through quick beats (travel, pickup, celebration). The humor depends on closure: we fill in all the unshown effort between panels (driving, meeting, negotiating) and then laugh when the payoff is presented as disproportionately triumphant. Because the format is small and fast, the comic compresses time, making ordinary actions feel like a mini-quest. The platform’s bite-sized pacing becomes part of the meaning: a big emotional reaction is squeezed into a tiny narrative container---exactly how social feeds often turn daily life into shareable highlights.

% ---------- Optional Figure Placeholder (insert your Catana 4-panel screenshot) ----------
% \begin{figure}[!t]
%     \centering
%     \includegraphics[width=\linewidth]{catana-placeholder.png}
%     \caption{CatanaComics: four-panel compression of time and effort into a quick payoff (insert your screenshot).}
%     \label{fig:catana}
% \end{figure}
% -------------------------------------------------------------------------------

\textit{Strange Planet} pushes scale in a different direction: it inflates the mundane into the cosmic through literal, scientific-sounding language \cite{CourseWeek12,Pyle2019}. Across the Week 12 examples, the aliens describe familiar experiences (``education,'' ``habitual examination,'' ``hypothetical futures'') in a way that reframes them as strange systems \cite{CourseWeek12}. Visually, the panels remain simple and evenly paced, with consistent framing that creates a calm rhythm. The scale shift happens mostly in the words. Because the drawings are minimal and the language is oddly grand, the reader performs closure by translating: we realize what human situation is being described, and the gap between the alien phrasing and the ordinary action becomes the punchline. This is an interdependent word--image strategy: neither the picture nor the text alone fully delivers the joke; meaning arrives when the reader fuses them. That fusion also produces a ``scale effect'': everyday life becomes both tiny (seen from an outsider’s distance) and enormous (re-described with cosmic seriousness).

% ---------- Optional Figure Placeholder (insert your Strange Planet 4-panel screenshot) ----------
% \begin{figure}[!t]
%     \centering
%     \includegraphics[width=\linewidth]{strangeplanet-placeholder.png}
%     \caption{\textit{Strange Planet}: ordinary scenes reframed through cosmic, literal language (insert your screenshot).}
%     \label{fig:strangeplanet}
% \end{figure}
% -------------------------------------------------------------------------------

\section{Comparison and Conclusion}
Comparing these works shows that scale is not just about epic vs.\ small art styles---it is about how comics structure perception. \textit{Watchmen} uses controlled pacing, symbolic zooming, and strategic gaps to make individual moments echo with political magnitude \cite{MooreGibbonsWatchmen}. The social media comics use compression, quick closures, and punchline timing to make minor experiences feel emotionally oversized---or to make human habits feel cosmically strange \cite{CourseWeek12,CatanaComics,Pyle2019}. In both cases, scale is produced by what comics do best: shaping time through sequence, shaping meaning through the gutter, and shaping emphasis through the interaction of words and images \cite{McCloud1993}. Whether the story is about conspiratorial power or a \$7 marketplace victory, comics can make the small feel huge---and remind us that ``huge'' is often a matter of framing.

\begin{thebibliography}{00}

\bibitem{McCloud1993}
S.~McCloud, \textit{Understanding Comics: The Invisible Art}. New York, NY, USA: HarperPerennial, 1993.

\bibitem{MooreGibbonsWatchmen}
A.~Moore and D.~Gibbons, \textit{Watchmen}. New York, NY, USA: DC Comics, 1986--1987. (Chapter 1 referenced).

\bibitem{CourseWeek12}
LSO~312 Course Materials, ``Week 12 Social Media Posts'' (CatanaComics and \textit{Strange Planet} examples), 2025. (Course PDF).

\bibitem{CatanaComics}
C.~Chetwynd, ``CatanaComics'' (webcomic/Instagram series), accessed via course Week 12 examples.

\bibitem{Pyle2019}
N.~W.~Pyle, \textit{Strange Planet}. New York, NY, USA: William Morrow, 2019. (Examples accessed via course Week 12 materials).

\end{thebibliography}

\end{document}
